% The calculus in this code — Section 1: the radial operator
% Compile with: xelatex calculus_in_this_code.tex
% Source of record: docs/guides/calculus_in_this_code.md

\documentclass[11pt]{article}

\usepackage{fontspec}
\setmainfont{Latin Modern Roman}
\setsansfont{Latin Modern Sans}
\setmonofont{DejaVu Sans Mono}[Scale=0.82]

\usepackage{unicode-math}
\setmathfont{Latin Modern Math}

\usepackage[a4paper,margin=1.05in,top=1.1in,bottom=1.15in]{geometry}
\usepackage{fancyvrb}
\usepackage{xcolor}
\usepackage{booktabs}
\usepackage{longtable}
\usepackage{array}
\usepackage{titlesec}
\usepackage{fancyhdr}
\usepackage[most]{tcolorbox}
\usepackage{enumitem}
\usepackage{microtype}
\usepackage{hyperref}

\definecolor{rule}{HTML}{9AA0A6}
\definecolor{panel}{HTML}{F4F5F7}
\definecolor{panelrule}{HTML}{C9CDD3}
\definecolor{codebg}{HTML}{F7F7F5}
\definecolor{coderule}{HTML}{DCDCD8}
\definecolor{heading}{HTML}{1A1A1A}

\hypersetup{
  colorlinks=true, linkcolor=black, urlcolor=black,
  pdftitle={The calculus in this code — Section 1: the radial operator},
  pdfsubject={What the radiodialysis solver actually integrates}
}

\titleformat{\section}
  {\normalfont\large\bfseries\color{heading}}{\thesection.}{0.6em}{}
\titlespacing*{\section}{0pt}{1.9\baselineskip}{0.7\baselineskip}
\titleformat{\subsection}
  {\normalfont\normalsize\bfseries\color{heading}}{}{0em}{}
\titlespacing*{\subsection}{0pt}{1.2\baselineskip}{0.4\baselineskip}

\pagestyle{fancy}
\fancyhf{}
\renewcommand{\headrulewidth}{0.4pt}
\renewcommand{\headrule}{\color{rule}\hrule height 0.4pt}
\fancyhead[L]{\footnotesize\color{rule}The calculus in this code}
\fancyhead[R]{\footnotesize\color{rule}Section 1 — the radial operator}
\fancyfoot[C]{\footnotesize\color{rule}\thepage}

\setlength{\parskip}{0.55\baselineskip}
\setlength{\parindent}{0pt}

% Code blocks: raw Unicode goes straight to DejaVu Sans Mono.
\DefineVerbatimEnvironment{code}{Verbatim}
  {fontsize=\small, xleftmargin=1.1em, xrightmargin=0.6em,
   frame=leftline, framerule=1.4pt, rulecolor=\color{coderule},
   framesep=0.85em, baselinestretch=1.02}

\newtcolorbox{panelbox}{
  colback=panel, colframe=panelrule, boxrule=0.4pt,
  arc=1.5pt, left=1.05em, right=1.05em, top=0.9em, bottom=0.9em,
  breakable,
  before upper={\setlength{\parskip}{0.55\baselineskip}}
}

% Check tables: fragments are code, set small and allowed to break.
\newcolumntype{C}{>{\raggedright\arraybackslash}p{0.285\textwidth}}
\newcolumntype{S}{>{\raggedright\arraybackslash\footnotesize\ttfamily}p{0.275\textwidth}}
\newcolumntype{F}{>{\raggedright\arraybackslash\footnotesize\ttfamily}p{0.315\textwidth}}

\newcommand{\us}{\_\allowbreak}
\newcommand{\checkhead}{%
  \toprule
  \textbf{\normalsize claim} & \textbf{\normalsize source} & \textbf{\normalsize must contain} \\
  \midrule
}

\begin{document}

\begin{center}
{\LARGE\bfseries The calculus in this code}\\[0.45em]
{\large Section 1 — the radial operator}\\[1.1em]
{\footnotesize\color{rule}%
Derived from \texttt{docs/guides/calculus\_in\_this\_code.md}, which is the checked source.\\
This rendering is a derived artifact and can go stale independently of it.}
\end{center}

\vspace{0.6\baselineskip}

\begin{panelbox}
One line of the solver carries the whole geometry:

\begin{code}
∂c/∂t = (1/r) ∂/∂r (r D_eff ∂c/∂r)
\end{code}

Two $r$s, and neither is a coordinate trick. The one inside is flux times area. The one
outside is division by volume. Everything in Part 1 follows from a shell between $r$ and
$r + \mathrm{d}r$ whose two faces have different areas, and you can get there without ever
writing down a coordinate transformation.

The other four parts are the places where that operator meets the code: why the axis needs
no boundary condition, why this solver special-cases it anyway, how a Robin wall becomes one
line of ghost-node algebra, and why forty coupled ODEs go to LSODA rather than a fixed step.

\textbf{Floor:} the chain rule, the product rule, and what a partial derivative is.
\textbf{Ceiling:} ordinary differential equations as objects an integrator solves.

Not here: general curvilinear coordinates, the divergence theorem in 3D, stability proofs,
and the sorption terms beyond noting that they set a second timescale.

Each part ends with a \emph{Checks against the source} table. Every row carries a
\texttt{file:line} and a fragment that line must contain, and
\texttt{calibration/tests/test\_guide\_citations.py} verifies all of them — a stale citation
fails the suite, and the failure says whether the code moved or the code changed.

\textbf{Claims in the prose but not in a table are not mechanically verified.} Three here:
the reading of \texttt{bc\_coef} as $P_{\text{eff}}$ versus $k_L$, the sorption rates in
Part 5, and the stiffness characterisation in Part 5. All three are sourced in the text.
None is pinned by a test. A guide implying uniform coverage would be making exactly the
claim its own tables exist to discipline.

\texttt{README.md}'s \emph{Mathematical framework} states the project's \emph{target
formalism}, under its own disclaimer that what the programs integrate differs. This is the
other side: what the code integrates. It cross-references those equations rather than
restating them, so the two cannot drift into disagreeing.
\end{panelbox}

\section{Where the \texorpdfstring{$r$}{r} comes from}

The solver's header states the equation it integrates:

\begin{code}
∂c/∂t = (1/r) ∂/∂r (r D_eff ∂c/∂r)
\end{code}

Most courses introduce this as ``the Laplacian in cylindrical coordinates'' and hand you a
table. That presentation makes the $r$ look like bookkeeping from a change of variables. It
isn't. It is a statement about area, and you can derive it without ever mentioning
coordinates.

\textbf{Take a shell.} Consider the region between radius $r$ and radius $r + \mathrm{d}r$,
one unit tall. Two facts about it:

\begin{itemize}[leftmargin=1.4em, itemsep=0.15em, topsep=0.3em]
\item its \textbf{volume} is $2\pi r\,\mathrm{d}r$ (circumference times thickness), and
\item its two \textbf{faces} have areas $2\pi r$ and $2\pi(r + \mathrm{d}r)$ —
      \emph{different areas}.
\end{itemize}

That difference is the whole story. In a slab, the two faces of a slice have the same area,
and everything below collapses to $\partial^2 c/\partial x^2$.

\textbf{Write conservation.} The amount of solute in the shell changes at the rate it enters
minus the rate it leaves. Rate of flow through a face is flux times area, so with $J(r)$ the
flux (amount per area per time, positive outward):
\[
\frac{\mathrm{d}}{\mathrm{d}t}\bigl[\,c \cdot 2\pi r\,\mathrm{d}r\,\bigr]
\;=\; J(r)\cdot 2\pi r \;-\; J(r+\mathrm{d}r)\cdot 2\pi(r+\mathrm{d}r)
\]

Divide by the volume $2\pi r\,\mathrm{d}r$:
\[
\frac{\partial c}{\partial t}
\;=\; -\,\frac{(r+\mathrm{d}r)J(r+\mathrm{d}r) - r\,J(r)}{r\,\mathrm{d}r}
\]

The numerator is the change in the product $rJ$ across the shell. Let
$\mathrm{d}r \to 0$ and it is by definition the derivative of that product:
\[
\frac{\partial c}{\partial t} \;=\; -\frac{1}{r}\frac{\partial (rJ)}{\partial r}
\]

\textbf{Now put Fick's law in.} Flux runs down the concentration gradient,
$J = -D\,\partial c/\partial r$. Substituting:
\[
\frac{\partial c}{\partial t}
\;=\; \frac{1}{r}\frac{\partial}{\partial r}\!\left( r D \frac{\partial c}{\partial r} \right)
\]

which is the line in the header, derived from a conservation statement and two areas.

\textbf{Read the two $r$s separately, because they are different things.} The $r$
\emph{inside} the derivative is flux $\times$ area — it is there because a face's area grows
with radius. The $1/r$ \emph{outside} is $\div$ volume — it is there because you asked for a
concentration rather than an amount. Neither is a coordinate artifact.

The physical content is real. Flux converging on a shrinking face concentrates;
flux spreading onto a growing face dilutes. A slab has neither effect, which is why its
operator has no such term. If you expand the product for constant $D$:
\[
\frac{\partial c}{\partial t}
\;=\; D\left[\, \frac{\partial^2 c}{\partial r^2}
  + \frac{1}{r}\frac{\partial c}{\partial r} \,\right]
\]

the first term is ordinary diffusion and \textbf{the second is purely geometric} — it is the
dilution-onto-a-larger-face effect, and it is proportional to the gradient rather than to its
curvature. Note that it carries a $1/r$, which is where Part 2 begins.

\subsection*{Checks against the source}

\begin{longtable}{CSF}
\checkhead
\endhead
the solver states the cylindrical operator
  & biofilms\us radiodialysis.R:9
  & (1/r) ∂/∂r (r D\_eff ∂c/∂r) \\[0.35em]
face weights are area-over-volume at the half-indices
  & biofilms\us radiodialysis.R:223
  & w\_plus  <- (r\_grid + 0.5 * dr) / r\_grid \\[0.35em]
and the inward face likewise
  & biofilms\us radiodialysis.R:224
  & w\_minus <- (r\_grid - 0.5 * dr) / r\_grid \\[0.35em]
geometry enters the stencil only through those weights
  & biofilms\us radiodialysis.R:120
  & (w\_plus[i]  * (c\_vec[i + 1] - c\_vec[i]) - \\
\bottomrule
\end{longtable}

\section{Why the axis isn't a boundary}

That $(1/r)\,\partial c/\partial r$ term is singular at $r = 0$. The solution is not.

The resolution is symmetry. The model is radially symmetric — $c$ depends on
distance from the axis and on nothing else. A smooth function of radius alone must have
$\partial c/\partial r \to 0$ at $r = 0$. If it did not, the profile would have a cusp on the
axis, and a cusp in a diffusion profile means a source sitting exactly there. There is no
such source in this model.

So at the axis the singular term is $0/0$, and the limit exists. Apply L'Hôpital to the
quotient $(\partial c/\partial r)/r$ as $r \to 0$: differentiate numerator and denominator
separately, giving $(\partial^2 c/\partial r^2)/1$. The geometric term therefore contributes
exactly as much as the ordinary term, and the operator becomes
\[
\lim_{r\to 0} D\left[\frac{\partial^2 c}{\partial r^2}
 + \frac{1}{r}\frac{\partial c}{\partial r}\right]
 \;=\; 2D\,\frac{\partial^2 c}{\partial r^2}
\]

\textbf{The factor of 2 is not a fudge.} It is the geometric term, evaluated in the limit,
turning out to equal the ordinary term. The cleanest way to remember it: the factor equals
the number of spatial dimensions the radial coordinate is standing in for — 2 here, because
$r$ is the radius of a disc in the plane; it would be 3 for a sphere. On the axis, diffusion
is that many times as effective at flattening curvature as it is in a slab.

This is why the manuscript calls it symmetry rather than a wall. The condition at
$r = 0$ is a \emph{consequence of the geometry}, not a physical boundary anybody chose. There
is no membrane on the axis, nothing is sealed, and no flux is being blocked. The zero
gradient is what radial symmetry forces, and the ``zero-flux'' phrasing describes the
consequence rather than a mechanism.

\subsection*{Checks against the source}

\begin{longtable}{CSF}
\checkhead
\endhead
the axis limit is named as L'Hôpital
  & biofilms\us radiodialysis.R:106
  & L'Hôpital limit \\[0.35em]
and implemented with the factor of 2
  & biofilms\us radiodialysis.R:110
  & dc\_dt[1] <- D\_eff * 2.0 * (c\_vec[2] - c\_vec[1]) / dr\^{}2 \\[0.35em]
the manuscript calls it symmetry, not a wall
  & preprint/modeling\us radioresistance\us and\us radiotropic\us fitness.tex:780
  & Zero-flux symmetry is imposed at \$r = 0\$. \\
\bottomrule
\end{longtable}

\section{Three treatments, three reasons}

Here is the part where a description of the \emph{method} and a description of \emph{this
code} come apart, and the gap is worth being precise about.

\textbf{The idealised finite-volume story.} Finite difference discretises the
\textbf{operator} — it approximates $\partial^2 c/\partial r^2$ and
$(1/r)\,\partial c/\partial r$ separately with difference quotients, so it meets the $1/r$
head-on and has to special-case the singular node. Finite volume discretises the
\textbf{conservation law} instead: integrate over each cell and what remains is face fluxes
times face areas. You never differentiate $1/r$ at all.

On that telling, both hard parts vanish for free. The innermost cell's inner face sits at
radius zero, so its \textbf{area} is zero, and the axis condition enforces itself by
multiplying a flux by nothing. At the wall, a Robin condition is a statement about flux — and
flux is already the quantity the scheme carries — so you set the outermost face flux
directly, with no ghost node and no one-sided derivative.

\textbf{This solver does neither of those things, and the interior is the only place the
story holds.}

\emph{The interior} is genuinely flux-form. Geometry enters through two face weights and
nowhere else:

\begin{code}
w_plus  <- (r_grid + 0.5 * dr) / r_grid
w_minus <- (r_grid - 0.5 * dr) / r_grid
\end{code}

Those are face area over cell volume, exactly the ratio Part 1 derived, evaluated at the
half-indices $r \pm \mathrm{d}r/2$ where the faces are. The stencil then reads as (outward
face flux $-$ inward face flux), weighted, and no $1/r$ is ever differentiated.

\emph{The axis is special-cased.} The innermost node does not get a zero-area face; it gets
the L'Hôpital limit from Part 2 written in directly:

\begin{code}
dc_dt[1] <- D_eff * 2.0 * (c_vec[2] - c_vec[1]) / dr^2 + ...
\end{code}

That line is the one Part 2's table already pins, so it is not repeated below. What the table
here adds is the second half of the evidence: the face weights at that node are set to
\texttt{NA} on purpose, so any future code that reaches for a metric weight there fails
loudly instead of silently using one.

\emph{The wall uses a ghost node.} Not a directly-set face flux — the eliminated-ghost
algebra that the idealised account says finite volume avoids:

\begin{code}
c_ghost <- c_vec[Nr - 1] -
  2.0 * dr * bc_coef * (c_vec[Nr] - c_ext) / D_eff
\end{code}

So the code is a hybrid, and the contrast is what makes it worth explaining. Flux-form in the
interior, L'Hôpital at the axis, ghost-node elimination at the wall — three treatments and
three reasons. The interior needs no special case \emph{because} it is flux-form; the two
boundaries need one \emph{because} the scheme is not carried through to the faces there. The
idealised account is correct about the method and it explains why the middle looks the way it
does. It is not a description of the two ends.

\textbf{None of that contradicts the manuscript.} \S5 describes the axis limit and the Robin
condition as represented explicitly in the semi-discrete operator, and both are: the
L'Hôpital limit and the ghost elimination are written out inside the right-hand side rather
than imposed on it from outside. ``Explicitly in the operator'' and ``through the face
structure'' are different claims, and only the second is the one this part is distinguishing
from.

\subsection*{Checks against the source}

\begin{longtable}{CSF}
\checkhead
\endhead
the axis node's face weights are deliberately unusable
  & biofilms\us radiodialysis.R:226
  & w\_plus[1] <- NA\_real\_; w\_minus[1] <- NA\_real\_ \\[0.35em]
the wall uses a ghost node, not a set face flux
  & biofilms\us radiodialysis.R:132
  & c\_ghost <- c\_vec[Nr - 1] - \\[0.35em]
the scheme is named finite-volume method of lines
  & biofilms\us radiodialysis.R:32
  & finite-volume method of lines \\[0.35em]
the manuscript describes both ends as handled in the operator
  & preprint/modeling\us radioresistance\us and\us radiotropic\us fitness.tex:1044
  & represented explicitly in the semi-discrete operator \\
\bottomrule
\end{longtable}

\section{The Robin condition}

At the outer wall the solver imposes:
\[
-D_{\text{eff}}\,\frac{\partial c}{\partial r}\bigg|_{r=R}
 \;=\; P_{\text{eff}}(t)\,\bigl(c(R,t) - c_{\text{ext}}\bigr)
\]

Read it as a flux balance. The left side is the diffusive flux arriving at the wall from
inside. The right side is what the membrane will pass, proportional to the concentration
difference across it, with $P_{\text{eff}}$ the constant of proportionality — a permeability,
in units of velocity.

It interpolates between the two conditions you already know.

\begin{itemize}[leftmargin=1.4em, itemsep=0.2em, topsep=0.3em]
\item $P_{\text{eff}} \to \infty$: the right side can only stay finite if
      $c(R) \to c_{\text{ext}}$. The membrane is invisible; this is a Dirichlet condition.
\item $P_{\text{eff}} \to 0$: the right side vanishes, so
      $\partial c/\partial r|_R = 0$. Nothing crosses; this is a sealed wall, a Neumann
      condition.
\end{itemize}

The membrane is the only thing in this model that is neither, and that is exactly why the
condition has to be Robin rather than one of the two simpler kinds.

Getting from that statement to the code is one substitution. Introduce a ghost node
$c[N_r{+}1]$ just outside the domain and approximate the wall derivative with a centred
difference, which is second-order:
\[
\frac{\partial c}{\partial r}\bigg|_R \approx
 \frac{c[N_r{+}1] - c[N_r{-}1]}{2\,\mathrm{d}r}
\]

Put that into the Robin condition and solve for the ghost:
\begin{align*}
-D\,\frac{c[N_r{+}1] - c[N_r{-}1]}{2\,\mathrm{d}r}
 &= \texttt{bc\_coef}\,\bigl(c[N_r] - c_{\text{ext}}\bigr) \\[0.25em]
c[N_r{+}1] &= c[N_r{-}1] - 2\,\mathrm{d}r\;\texttt{bc\_coef}\,
 \bigl(c[N_r] - c_{\text{ext}}\bigr) / D
\end{align*}

which is the line in the solver. The ghost is then substituted into the ordinary interior
stencil, so it never appears in the state vector — it exists only inside the arithmetic for
the last node.

One naming caution, and the code is emphatic about it. \texttt{bc\_coef} is
$P_{\text{eff}}$ in the cylindrical geometry and $k_L$, an external liquid-film mass-transfer
coefficient, in the slab. They are different physical quantities, and the slab's must not be
reported as a permeability. The stencil is written to take whichever the producer declares.

\subsection*{Checks against the source}

\begin{longtable}{CSF}
\checkhead
\endhead
the condition as stated
  & biofilms\us radiodialysis.R:20
  & -D\_eff ∂c/∂r|\_\{r=R\} = P\_eff(t) · (c(R,t) − c\_ext) \\[0.35em]
the ghost substitution, second line
  & biofilms\us radiodialysis.R:133
  & 2.0 * dr * bc\_coef * (c\_vec[Nr] - c\_ext) / D\_eff \\
\bottomrule
\end{longtable}

\section{Method of lines, and why LSODA}

\textbf{The construction.} Discretise space and leave time alone. Forty cells, each with a
concentration, gives forty coupled ordinary differential equations — plus forty more for the
sorbed phase and one for membrane integrity, packed into a single state vector of length
$2N_r + 1$. That is the method of lines, and the point of it is that once space is
discretised you are holding an ODE system, and ODE integrators are a solved problem you can
hand it to.

Why not step it forward yourself? An explicit step is limited by stability, not
by accuracy. The diffusion operator's eigenvalues scale as $D/\Delta r^2$, and an explicit
Euler step is stable only while $|1 + \Delta t\,\lambda| \le 1$ for every eigenvalue
$\lambda$ — which caps the step at roughly $\Delta r^2/2D$. Halve the grid spacing and the
allowed step drops by four.

That bound is not hypothetical here: the Julia port's explicit stepper substeps against
exactly it, with a safety factor:

\begin{code}
dt_stable = 0.4 * dr^2 / (2.0 * rd.params.D_eff)
\end{code}

That is only one of the timescales. Sorption runs at its own rate — an uptake of
$U = 0.056\ \mathrm{s^{-1}}$ against a relaxation of
$1/(k_{\text{des}} + k_{\text{loss}}) = 166.7\ \mathrm{s}$. When the fastest process in a
system forces a step far shorter than the slowest process needs for accuracy, the system is
\textbf{stiff}, and that is the whole of the definition.

\textbf{Whether \emph{this} system is stiff is not computed anywhere in this repository, and
this guide does not assert it.} Stiffness is a ratio, and no artifact here reports one:
\texttt{analysis/verify\_radiodialysis\_stability.py} computes the explicit-Euler bound and
the conservative real-axis limit, but not a stiffness ratio against the sorption timescale.
The two rates quoted above are the ingredients such a ratio would be built from, and they are
themselves among the prose claims no test pins. So the honest statement is that LSODA is a
defensible default for a system with two visibly separated timescales, not a measured
requirement — and closing that gap means shipping the ratio from a producer, in the shape the
rest of this repository already uses.

That is what LSODA is for. It monitors the solution and switches between an Adams
method (cheap, for the non-stiff stretches) and BDF (implicit, stable at large steps, for the
stiff ones), choosing its own step size to meet the tolerances it is given rather than a
stability limit you computed by hand.

\textbf{A caveat this repository is careful about, and so is this guide.} The explicit bound
above is analysed in \texttt{analysis/verify\_radiodialysis\_stability.py}, which separates
three claims that are easy to conflate — continuous stability, semi-discrete stability, and
explicit-Euler stability — and computes the conservative real-axis limit directly. Its
conclusion about the Julia path is that the explicit stepper is stable at
$\mathrm{d}t = 0.5$ only because $r$ is in lattice units while $D_{\text{eff}}$ is in
$\mathrm{cm^2\,s^{-1}}$, which is the documented units error behind
\texttt{RADIODIALYSIS: BLOCKED}. \textbf{A stability margin that comes from a units mismatch
is not evidence the step is safe}, and none of the reasoning in this part should be read as
saying the coupled path is sound.

\subsection*{Checks against the source}

\begin{longtable}{CSF}
\checkhead
\endhead
forty cells by default
  & biofilms\us radiodialysis.R:230
  & default\_parms <- function(Nr = 40, R = 1.0) \\[0.35em]
the state vector packs c, s and m
  & biofilms\us radiodialysis.R:50
  & y[1 .. Nr]        = c\_i \\[0.35em]
LSODA is the integrator
  & biofilms\us radiodialysis.R:400
  & method = "lsoda" \\[0.35em]
the Julia port substeps against the diffusion bound
  & biofilms\us potts.jl:1453
  & dt\_stable = 0.4 * dr\^{}2 / (2.0 * rd.params.D\_eff) \\[0.35em]
the explicit-Euler limit is computed, not asserted
  & analysis/verify\us radiodialysis\us stability.py:153
  & conservative\_real\_axis\_limit \\
\bottomrule
\end{longtable}

\vspace{1.2\baselineskip}
{\footnotesize\color{rule}\hrule\vspace{0.7em}
Every check-table row is verified by
\texttt{calibration/tests/test\_guide\_citations.py}: the file must exist, the line must
exist, and it must contain the fragment. A stale line number fails the suite, and the failure
message distinguishes ``the code moved'' from ``the code changed.'' Long paths are broken across lines at underscores; the breaks are typographic only.}

\end{document}
